\documentclass[12pt]{article}
\usepackage[T1]{fontenc}
\usepackage[utf8]{inputenc}
\usepackage{lmodern}
\usepackage{textcomp}
\usepackage{enumitem}
\usepackage{geometry}
\geometry{margin=1in}
\begin{document}
\begin{center}
Answer Key - Religious Studies - K-12 Level Assessment 16 \
\small (Answers are ordered exactly as in the question paper.)
\end{center}

\section*{Multiple Choice}
\begin{enumerate}[label=\arabic*.]
\item \textbf{Answer:} C
\\ \textit{Options:} A) Piety; B) Non-action; C) Heart-mind; D) World
\\ \textit{Note:} The term xin in Chinese philosophy and religious studies is commonly understood to denote the interconnectedness of the heart and mind, reflecting the idea that emotional and cognitive processes are closely linked.
\item \textbf{Answer:} D
\\ \textit{Options:} A) The Classic of Documents; B) The Lotus Sutra; C) The Flower Garland Sutra; D) The Classic of Changes
\\ \textit{Note:} The correct answer is "The Classic of Changes" because it is the English translation of the Chinese title "Yijing," which refers to an ancient text used for divination and philosophical insight in Chinese culture.
\item \textbf{Answer:} A
\\ \textit{Options:} A) 1935; B) 2000; C) 1965; D) 1975
\\ \textit{Note:} The current Dalai Lama, Tenzin Gyatso, was born on July 6, 1935, in Tibet, making this date the correct answer to the question.
\item \textbf{Answer:} B
\\ \textit{Options:} A) Presbyteries; B) Sees; C) Churches; D) Synods
\\ \textit{Note:} The term "sees" refers to the five principal centers of authority and governance in early Christianity, each associated with a specific geographic region and bishopric, which played a crucial role in the development of Christian doctrine and organization.
\item \textbf{Answer:} D
\\ \textit{Options:} A) Honen; B) Tanaka; C) Tokugawa; D) Meiji
\\ \textit{Note:} The Meiji government promoted Shinto as a state religion, emphasizing the divine status of the emperor and his connection to kami, which solidified a national cult centered around imperial authority and Japanese identity.
\end{enumerate}

\section*{True / False}
\begin{enumerate}[label=\arabic*.]
\setcounter{enumi}{5}
\item \textbf{Answer:} False
\\ \textit{Options:} A) True; B) False
\\ \textit{Note:} Religious traditions in China and Korea, such as Confucianism, Buddhism, and Daoism, primarily emphasize personal spirituality, ethical living, and harmony with the community and nature rather than the pursuit of power and influence over others.
\item \textbf{Answer:} False
\\ \textit{Options:} A) True; B) False
\\ \textit{Note:} Ordinary folk in the Han Dynasty primarily sought assistance from local deities and spirits rather than Laozi, who was associated with philosophical teachings rather than direct intervention in everyday hardships.
\item \textbf{Answer:} False
\\ \textit{Options:} A) True; B) False
\\ \textit{Note:} The term Gathas refers to a collection of hymns attributed to Zoroaster, which express spiritual and philosophical teachings rather than specific ethical laws governing behavior.
\item \textbf{Answer:} False
\\ \textit{Options:} A) True; B) False
\\ \textit{Note:} The Hebrew word "mashiach" actually means "anointed one" and refers to a messianic figure or leader in Jewish tradition, rather than simply signifying a prophet or religious teacher.
\item \textbf{Answer:} True
\\ \textit{Options:} A) True; B) False
\\ \textit{Note:} The title 'Mahatma' is derived from the Sanskrit words 'maha' meaning 'great' and 'atma' meaning 'soul,' thus translating to 'Great soul' in English.
\end{enumerate}

\section*{Long Form Response}
\begin{enumerate}[label=\arabic*.]
\setcounter{enumi}{10}
\item \textbf{Answer:} Hwanung's transformation into a human is significant in the Korean foundation myth as it symbolizes the connection between the divine and the earthly realm. According to the myth, Hwanung descended from heaven to establish a society and brought order and culture to the people. This act of becoming human reflects Korean cultural beliefs about the importance of harmony between the spiritual and physical worlds, as well as the value placed on leadership and guidance from divine sources. Moreover, Hwanung's role emphasizes the idea of nurturing civilization and the responsibility of leaders to promote the well-being of their communities. Ultimately, this transformation highlights the belief in the potential for divine qualities to manifest in human life, shaping the cultural identity of Korea.
\\ \textit{Note:} correct because it emphasizes Hwanung's transformation as a pivotal moment that illustrates the integration of divine authority in human affairs, reflecting key Korean cultural values of harmony, leadership, and the interplay between the spiritual and earthly realms.
\item \textbf{Answer:} Guru Nanak began preaching the message of the divine Name at the age of 30, marking a significant turning point in his life and the development of Sikhism. This age is important as it symbolizes a moment of spiritual awakening and commitment to sharing his insights about God and equality. By choosing to spread his teachings at this time, Guru Nanak laid the foundation for Sikhism, emphasizing devotion to one God, the importance of community service, and the rejection of caste distinctions. His teachings attracted followers and established a new religious path that encouraged individuals to seek a direct relationship with the divine. This pivotal moment not only shaped the core beliefs of Sikhism but also inspired a movement that promoted peace, tolerance, and social justice.
\\ \textit{Note:} correct because it emphasizes that Guru Nanak's decision to start preaching at the age of 30 represents a pivotal moment in his spiritual journey, which ultimately led to the establishment of Sikhism's core tenets of monotheism, equality, and community service.
\end{enumerate}

\vfill
\noindent\textit{End of Answer Key}
\end{document}